% ЗГЕНЕРОВАНО scripts/build_format_v2_chapter.py — не редагувати вручну.
% Джерело: books/ashtavakra-gita/poetychnyi/rozdil-17.md

\poemverse{17.1}{Хто вдоволений, \\ чиї чуття прозорі \\ і хто завжди сам — \\ той і справді радіє: \\ він уже має плід знання \\ і плід вправ у йозі.}
\poemverse{17.2}{Ніколи в цьому світі не мучиться \\ той, хто знає істину: \\ о, тим одним сповнене все — \\ усе коло світів.}
\poemverse{17.3}{Жодні чуттєві речі не тішать того, \\ кому втіха — у ньому самому, \\ як не тішить слона, що любить пагони салакі, \\ гірке листя німби.}
\poemverse{17.4}{Хто вже зазнаних утіх не тягне за собою, \\ як просякнутий їхнім запахом, \\ і не прагне незазнаних — \\ така людина в цьому світі рідкість.}
\poemverse{17.5}{У цьому світі видно і тих, хто голодний до нього, \\ і тих, хто прагне звільнення. \\ Але хто не прагне ні втіх, ні звільнення — \\ такий великий і рідкісний.}
\poemverse{17.6}{Для великодушної людини \\ немає відкинути й узяти — \\ ні щодо дгарми, артги, ками й мокші, \\ ні щодо життя й смерті.}
\poemverse{17.7}{Він не бажає, щоб світ загинув, \\ і не має відрази до того, що світ триває. \\ Тому блаженний сидить собі спокійно, \\ як завгодно, — \\ і живе з того, що має.}
\poemverse{17.8}{Цим знанням я досяг мети — \\ і з цим, уже розтанувши розумом, \\ він живе собі як завгодно: \\ дивиться, слухає, торкається, нюхає, їсть.}
\poemverse{17.9}{Порожній погляд, \\ марні рухи, \\ чуття ніби понівечені; \\ ні прагнення, ні відвернення — \\ бо океан сансари в ньому висох.}
\poemverse{17.10}{Він не пильнує і не спить, \\ не розплющує й не заплющує очей. \\ О диво: розум вільної людини \\ перебуває в якомусь зовсім іншому стані.}
\poemverse{17.11}{Усюди він у собі, \\ усюди його нутро чисте; \\ вільний від усіх звичних слідів, \\ звільнений сяє всюди.}
\poemverse{17.12}{Дивиться, слухає, торкається, \\ нюхає, їсть, бере, \\ говорить, ходить — \\ і вільний від бажаного й небажаного. \\ Він справді звільнений, цей великий душею.}
\poemverse{17.13}{Не гудить і не хвалить, \\ не тішиться й не гнівається, \\ не дає й не бере — \\ звільнений усюди без смаку до всього.}
\poemverse{17.14}{Побачивши жінку, повну пристрасті, \\ або смерть, що вже насувається, \\ він не збентежиться: \\ він у собі, він справді звільнений, \\ цей великий душею.}
\poemverse{17.15}{У щасті й у горі, \\ в чоловікові й у жінці, \\ в удачі й у біді — \\ для стійкого, хто всюди бачить однаково, \\ різниці немає.}
\poemverse{17.16}{Ні насильства, ні жалю, \\ ні зарозумілості, ні приниженості, \\ ні здивування, ні збурення — \\ немає цього в людині, \\ в якої скінчилося блукання.}
\poemverse{17.17}{Звільнений не ненавидить чуттєвих речей \\ і не жадібний до них; \\ його розум ні до чого не привʼязаний, \\ і він бере те, що прийшло, \\ і те, що не прийшло.}
\poemverse{17.18}{Розрізнень зібраний я чи ні, \\ корисно це чи шкідливо \\ порожній розум уже не знає: \\ він стоїть ніби в самотності.}
\poemverse{17.19}{Без моє і без я роблю, \\ з певністю, що немає нічого, — \\ усі його надії розтанули зсередини, \\ і він навіть діючи не діє.}
\poemverse{17.20}{Він вільний і від ясності розуму, \\ і від затьмарення, \\ і від сну, і від заціпеніння — \\ і досяг якогось невідомого стану, \\ цей чоловік із розталим розумом.}
